% Common macros for the online blueprint.

% THEOREMS
\newtheorem{theorem}{Theorem}[section]
\newtheorem*{theorem*}{Theorem}
\newtheorem{corollary}[theorem]{Corollary}
\newtheorem{hypothesis}[theorem]{Hypothesis}
\newtheorem{proposition}[theorem]{Proposition}
\newtheorem{lemma}[theorem]{Lemma}
\newtheorem{definition}[theorem]{Definition}
\newtheorem{conjecture}[theorem]{Conjecture}
\newtheorem{example}[theorem]{Example}
\newtheorem{question}[theorem]{Question}
\newtheorem{remark}[theorem]{Remark}

% NUMBER SYSTEMS
\newcommand{\Z}{\mathbb{Z}}
\newcommand{\Q}{\mathbb{Q}}
\newcommand{\R}{\mathbb{R}}
\newcommand{\C}{\mathbb{C}}
\newcommand{\N}{\mathbb{N}}
\newcommand{\F}{\mathbb{F}}

% NUMBER-THEORETIC OBJECTS
\newcommand{\OK}{\mathcal{O}_K}                  % Ring of integers of K
\newcommand{\Ocirc}{\mathcal{O}}                 % Generic ring of integers
\newcommand{\zetap}{\zeta_p}                     % Primitive p-th root of unity
\newcommand{\Qzetap}{\Q(\zeta_p)}                % p-th cyclotomic field
\newcommand{\Qzetapplus}{\Q(\zeta_p)^+}          % Maximal real subfield
\newcommand{\Cl}[1]{\mathrm{Cl}(#1)}             % Class group
\newcommand{\hp}{h(\Q(\zeta_p))}                 % Class number of Q(zeta_p)
\newcommand{\hplus}{h^+}                          % Plus part of class number
\newcommand{\hminus}{h^-}                         % Minus part of class number
\newcommand{\Gal}[1]{\mathrm{Gal}(#1)}            % Galois group
\newcommand{\Bk}{B}                               % Bernoulli number
\newcommand{\Bgen}[2]{B_{#1,#2}}                  % Generalised Bernoulli
\newcommand{\Norm}[1]{\mathrm{N}(#1)}             % Field norm

% L-functions and zeta
\newcommand{\Lfn}[2]{L(#1, #2)}                   % Generic L-value
\newcommand{\zetaK}[1]{\zeta_{#1}}                 % Dedekind zeta of K
\newcommand{\sinZeta}{\mathrm{sinZeta}}            % Hurwitz sine zeta
\newcommand{\zetaodd}{\zeta^{\mathrm{odd}}}        % Hurwitz odd zeta

% DELIMITERS
\newcommand{\set}[1]{\left\{ #1 \right\}}
\newcommand{\setof}[2]{\left\{ #1 \;\middle|\; #2 \right\}}
\newcommand{\abs}[1]{\left\lvert #1 \right\rvert}
\newcommand{\norm}[1]{\left\lVert #1 \right\rVert}
\newcommand{\ang}[1]{\left\langle #1 \right\rangle}

% MISC
\newcommand{\todo}[1]{{\color{red}{\textbf{TODO: #1}}}}
